\chapter{Reliability and Fault Tolerance}
\label{ch:reliability}

\section{Reliability in Distributed Systems}
\label{sec:reliability_overview}

A monitoring and anomaly detection framework must itself be reliable. A system that
crashes or produces stale results under load provides weaker guarantees than no
monitoring at all, since it creates a false sense of operational visibility. In the
context of this framework, reliability encompasses three concerns: the continuous
availability of the anomaly detection pipeline at the edge, the durability of the
federated learning coordination process, and the consistency of the persistent state
(causal graphs, model checkpoints, anomaly scores) across crashes and restarts.

Achieving reliability in a distributed setting requires deliberate design choices
around failure detection, redundancy, state persistence, and graceful degradation.

\section{CAP Theorem Implications}
\label{sec:reliability_cap}

The CAP theorem (introduced in Section~\ref{subsec:bg_cap}) constrains the guarantees
any distributed system can simultaneously offer under a network partition. This
framework prioritizes \emph{Availability} and \emph{Partition Tolerance} (AP),
accepting the possibility of transient consistency divergence rather than halting
operation during a partition event.

This choice is motivated by the nature of the workload:

\begin{itemize}
    \item \textbf{Federated Learning Rounds}: Rounds are inherently asynchronous and
    tolerant of partial participation. If a network partition isolates a subset of
    edge nodes, the coordinator proceeds with the available participants rather than
    blocking the round indefinitely. The resulting global model may be aggregated from
    fewer nodes but remains useful and is corrected in subsequent rounds.

    \item \textbf{Anomaly Detection at the Edge}: Operates on locally cached model
    weights and does not require a live connection to the coordinator. Edge nodes
    continue inference autonomously under partition, sacrificing global model
    freshness for continued local detection capability.

    \item \textbf{Root Cause Analysis}: Relies on the causal dependency graph, which
    is updated incrementally. Temporary staleness in the graph---caused by missed
    trace updates during a partition---is acceptable, as the graph converges once
    connectivity is restored.
\end{itemize}

The coordinator itself does not enforce distributed consensus on FL round state,
avoiding the latency and complexity of protocols such as Raft or Paxos \cite{ongaro2014raft, 
lamport2001paxos}. Instead, round state is persisted durably before each commit, 
enabling crash recovery without requiring agreement across replicas.

\section{Single Points of Failure}
\label{sec:reliability_spof}

\subsection{Identification}
\label{subsec:spof_identification}

The following components represent potential single points of failure (SPoFs) in the
framework:

\begin{itemize}
    \item \textbf{Central Coordinator}: The coordinator orchestrates all federated
    learning rounds and serves as the authoritative source of the global model. If it
    fails, FL rounds halt and edge nodes cannot receive updated weights until it
    recovers.

    \item \textbf{Trace Collector}: The component responsible for ingesting
    distributed traces and updating the causal dependency graph. A failure here
    prevents graph updates, causing root cause analysis to operate on a stale
    topology.

    \item \textbf{Persistent Storage}: The coordinator and edge nodes rely on a
    durable database for storing model checkpoints, round state, and the causal graph.
    Disk corruption or filesystem failures can result in loss of this state if
    redundancy mechanisms are not in place.
\end{itemize}

\subsection{Mitigation Strategies}
\label{subsec:spof_mitigation}

Each identified SPoF is addressed through a corresponding mitigation:

\begin{itemize}
    \item \textbf{Coordinator Redundancy}: The coordinator is designed as a stateless
    process whose authoritative state resides entirely in durable storage. This
    enables an active-passive failover setup, where a standby coordinator instance can
    take over by reading the last committed state from the database without requiring
    state transfer between processes.

    \item \textbf{Graceful Degradation at the Edge}: Edge nodes that cannot reach the
    coordinator after a configurable timeout switch to standalone inference mode,
    continuing anomaly detection with the last received global model. This ensures
    that a coordinator outage does not propagate into a detection outage.

    \item \textbf{Trace Collector Resilience}: Incoming trace data is buffered in a
    durable message queue before being consumed by the graph update process. A
    collector restart drains the buffer and catches up without data loss.
\end{itemize}

\section{Redundancy}
\label{sec:reliability_redundancy}

\subsection{Data Redundancy}
\label{subsec:redundancy_data}

The following persistent data assets are protected through redundancy mechanisms:

\begin{itemize}
    \item \textbf{Global Model Checkpoints}: The coordinator maintains a versioned
    model registry, retaining the last $N$ global model snapshots. On recovery, the
    most recent valid checkpoint is loaded and broadcast to edge nodes at the start of
    the next round.

    \item \textbf{Causal Dependency Graph}: The graph is persisted incrementally in
    durable storage with Write-Ahead Logging (WAL) enabled. WAL mode ensures that 
    in-progress writes do not corrupt the on-disk state during a crash, and readers 
    are never blocked by concurrent writers.

    \item \textbf{FL Round State}: Partial aggregation state (e.g., participating node list,
    received gradient updates, round sequence number) is written to durable storage
    before any round transition is committed. An incomplete round detected on restart
    is discarded and re-initiated from the last committed global model.
\end{itemize}

\section{Fault Tolerance Mechanisms}
\label{sec:fault_tolerance_mechanisms}

\subsection{Crash Detection}
\label{subsec:ft_crash_recovery}

The coordinator maintains a heartbeat-based health monitoring channel with all
registered edge nodes. Each node emits a periodic heartbeat over a lightweight UDP
channel at a fixed interval. If the coordinator does not receive a heartbeat from 
a node within a timeout window, the node is marked as offline and excluded from the 
current FL round's aggregation. Excluded nodes are not penalized; they are automatically 
re-enrolled when their heartbeat resumes.

Edge nodes perform symmetric health monitoring of the coordinator. A failed
coordinator connection triggers the standalone inference fallback described in
Section~\ref{subsec:spof_mitigation}. Upon reconnection, the node re-registers with
the coordinator and is enrolled in the next available FL round.

\subsection{Edge Node Profiles}
\label{subsec:ft_node_profiles}

Edge nodes operate under one of two profiles that govern their fault tolerance
behavior:

\begin{itemize}
    \item \textbf{Lightweight Profile}: Designed for resource-constrained devices
    (e.g., IoT sensors, minimal containers). The node is fully stateless: it holds no
    persistent local storage, joins an FL round, trains on local data, submits
    gradients, and disconnects. On crash, it simply re-registers with the coordinator
    at startup. Recovery is immediate but the node cannot resume an interrupted
    training epoch.

    \item \textbf{Standard Profile}: Designed for persistent edge appliances and
    servers. The node caches the last received global model checkpoint in local 
    durable storage. On crash recovery, it can resume an interrupted training epoch 
    from the cached model without waiting for the coordinator to broadcast the next 
    round, reducing recovery latency and avoiding a missed participation slot.
\end{itemize}
